\section{Revised Question and Scope}\label{sec:question}
Following the feedback, the label-free question below is now our main question, and scaling and annotation come after it. scMINER clusters cells with MICA, which scores a pair of cells by the mutual information between their binned expression profiles \cite{pan2025scminer}. Its benchmark uses the labels to set up every method. Each method either gets the true number of cell types, or keeps the run that hits that number with the best adjusted Rand index (ARI) against the reference labels. Someone clustering new data has neither the count nor the labels, so we ask:

\textbf{Main question.} When each method must choose its own final clustering without the reference labels, does MICA give partitions closer to the known cell types than principal components and Leiden clustering \cite{traag2019louvain} on the same genes?

Both methods get the same label-free rules, two so far. One is each method's default: Scanpy's Leiden resolution 1.0, and MICA's 1.822, scored at 1.8, the nearest value on its sweep. MICA's default sits in its code with no stated origin, so it may not be fully label-free. The other keeps the setting with the highest mean silhouette in the method's own embedding (the space MICA clusters in, and the principal components for Leiden), over $k$ from 2 to 20 or over the method's resolution sweep. The labels come in afterwards, only to compute ARI and the adjusted mutual information (AMI). Results where the labels choose the setting, as in the paper, are reported separately as a label-informed reference.

\textbf{Secondary question.} Can MICA's neighbor graph be built much faster without losing the neighbors that matter? We test it on subsets where brute force gives the exact answer, and use them to estimate the cost of the full atlas.

\textbf{Later, if time allows.} scMINER also infers the activity of regulators, such as transcription factors, from its gene networks. Does that activity add supported evidence beyond marker genes when a retrieval-augmented language model names clusters?

The anchor datasets are the two the feedback named, in the processed form the scMINER authors released, which we use throughout. PBMC14K has 13,605 blood cells in 7 types sorted with antibodies before sequencing, so its labels do not come from clustering \cite{zheng2017massively}. Zeisel has 2,903 of the 3,005 mouse brain cells in the original study, in 7 classes whose labels did come from the original authors' clustering \cite{zeisel2015cell}. We add eight more of their benchmark sets for breadth, and a 100,000-cell subset of the 2.3 million 10x v3 cells in the Allen whole mouse brain atlas \cite{yao2023allen} for the scaling test.
